\subsection{Finite checks and ordinary proofs have different jobs}
The graph verifier imports neither graph producer. It checks the literal source
hash, every topological constant/scale constructor and source-row pullback,
the complete logical provenance of every factor, both graph simplicity
conditions, every free spanning tree, and the alias and unit constraints.
It then independently expands the sum of squares for $2P$, divides each even
coefficient by two, and compares every coefficient with the delivered polynomial.
It checks coordinate liveness, total and individual degrees, coefficient height,
the constant term, and all reported counts. Eight concrete corruptions of a
profile must be rejected. The check remains active with Python optimization.

The direct arithmetic checks compare the unchanged inherited source prefix
with its authenticated parent, verify every new semantic row as an exact
polynomial identity, replay the affine export and nonnegative or integer
restoration certificates, and expand the complete sum of squared residuals.
In particular the domain of a removed coordinate is part of a restoration
certificate; a rational identity alone does not justify an integer or
nonnegative-coordinate elimination. The signed cap and sequential recovery
arguments additionally receive ordinary mathematical review.

The eleven-variable branch checks all six actual fixed coefficient payloads,
the literal binary arithmetic graphs, the entire three-radical norm identity,
all new local algebraic identities, and the highest-degree forms at the
authentic coefficient specialization. Its degree statements concern the
whole polynomial, including cancellation, and not only values on its zero set.

The finite checks do not search for a zero, establish that no zero exists,
or use a finite sample as proof of an unbounded Pell or consistency theorem.
The mathematical arguments establishing those implications are supplied in
the report and retained source texts. Recorded review notes state their scope.

\subsection{Reproduction}
The supported entry point is
\begin{quote}\ttfamily
python3 -B verify.py
\end{quote}
from the extracted archive directory. Python's standard library suffices for
the supported new checkers. The runner authenticates the static file manifest
and invokes the component checks with their documented paths. It preserves
the delivered equations. Component receipts record both normal and optimized
checks; the final aggregate receipt records exactly what the release run did.
The component README files document deterministic rebuilds into temporary
directories and their complete byte comparisons. A successful algebraic check
is not labeled a proof-assistant check.

The manuscript builds with a standard TeX Live installation:
\begin{quote}\ttfamily
cd report\\
latexmk -pdf -interaction=nonstopmode -halt-on-error main.tex
\end{quote}
TeX's PDF metadata can vary between builds. Deterministic byte-reproduction
claims concern the polynomial, circuit, and finite-certificate files. The
distributed PDF is itself covered by the release integrity manifest and was
rendered and visually inspected.

\subsection{How to read the literal files}
For a sparse polynomial, a term is written as
\file{[coefficient, [indices]]}. Thus $[-3,[7,7]]$ denotes $-3z_7^2$,
and $[1,[]]$ denotes the constant one. Coordinates are zero-based.
In a residual-system file, the exact final polynomial is the sum of the squares
of every residual, in its printed variable order. An expanded polynomial is
supplied as an additional literal representation. The graph files instead list
their final height-one polynomial directly and carry its half-sum-of-squares
certificate. The eleven-variable files use literal constant, variable, and
binary add/subtract/multiply nodes. A named intermediate node adds no unknown.

The manifest and source hashes detect accidental changes and bind the exact
released data. They do not establish a mathematical theorem by themselves.
